\pagestyle{plain}

In this section, we recall the essential definitions that will serve as our main framework throughout this article. 

\subsection{Vocabulary}
\label{ss:Vocab}

A \new{quiver} is a quadruple $Q = (Q_0,Q_1,s,t)$ where:
\begin{enumerate}[label=$\bullet$, itemsep=1mm]
    \item $Q_0$ is a set called \emph{vertex set};
    \item $Q_1$ is a set called \emph{arrow set}; and,
    \item $s,t : Q_0 \longrightarrow Q_1$ are functions called \emph{source} and \emph{target functions}.
\end{enumerate}
Such a quiver $Q$ is said to be \new{finite} whenever $Q_0$ and $Q_1$ are finite sets. In the following, we assume that all our quivers are finite.

A \new{path} of $Q$ is either a formal variable $e_q$ for some $q \in Q_0$ called \emph{lazy path at $q$}, or a finite nonempty word $p = \alpha_k \cdots \alpha_1$ on $Q_0$ such that, for all $i \in \{1,\ldots,k-1\}$, we have $s(\alpha_{i+1}) = t(\alpha_i)$. Given such a path $p$, we define:
\begin{enumerate}[label=$\bullet$, itemsep=1mm]
    \item its \emph{source} $s(p)$ to be $s(\alpha_1)$ (we set $s(e_q) = q$);
    \item its \emph{target} $t(p)$ to be $t(\alpha_k)$ (we set $t(e_q) = q$); and,
    \item its \textit{length} $\ell(p)$ to be $k$ (we set $\ell(e_q) = 0$).
\end{enumerate}
For any arrow $\alpha \in Q_1$, we consider its formal inverse $\alpha^{-1}$ such that $s(\alpha^{-1}) = t(\alpha)$ and $t(\alpha^{-1}) = s(\alpha)$. We denote by $\overline{Q_1}$ the set of formal inverses of arrows in $Q_1$. A \new{walk} of $Q$ is either a lazy path or a finite nonempty word $w = \beta_k \cdots \beta_1$ on $Q_1 \sqcup \overline{Q_1}$ such that, for all $i \in \{1,\ldots,k-1\}$, we have $s(\beta_{i+1}) = t(\beta_i)$. We extend the source, target, and length functions to walks of $Q$ similarly to the case of paths of $Q$. A quiver $Q$ is said to be \new{connected} whenever for any pair $(a,b) \in Q_0$, there exists a walk $w$ of $Q$ such that $s(w) = a$ and $t(w)=b$. From now on, we assume that all our quivers are connected.

\subsection{Representations}
\label{ss:Reps}

Let $\mathbb{K}$ be a field. For geometric purposes, we assume that $\mathbb{K}$ is algebraically closed. The \new{path algebra} of $Q$, denoted by $\mathbb{K}Q$, is the $\mathbb{K}$-vector space generated as a basis by the set of all the paths of $Q$, together with a multiplication which acts as concatenation on the paths: meaning \[p_2 \cdot p_1 = \begin{cases}
    p_2p_1 & \text{if } s(p_2) = t(p_1) \text{; and,} \\
    0 & \text{otherwise.}
\end{cases}\] For any $m \in \mathbb{N}$, we denote by $\mathbb{K}Q_{\geqslant m}$ the (bi-sided) ideal of $\mathbb{K}Q$ generated by all the paths of length greater than or equal to $m$. An ideal $I \subseteq \mathbb{K}Q$ is said to be \new{admissible} whenever there exists an integer $m \geqslant 2$ such that $\mathbb{K}Q_{\geqslant m} \subseteq  I \subseteq \mathbb{K} Q_{\geqslant 2}$. A \new{relation} on $Q$ is a (finite) $\mathbb{K}$-linear combination of paths of length at least two, and having a common source and a common sink. For any subset of relations $R$ on $Q$, we denote by $\langle R \rangle$ the ideal of $\mathbb{K}Q$ generated by $R$. Call \new{bounded quiver} any pair $(Q,R)$ where $Q$ is a quiver and $R$ is a set of relations on $Q$.

Given a bounded quiver $(Q,R)$, we define a \new{representation} of $Q$ (over $\mathbb{K}$) as a pair $E=((E_q)_{q \in Q_0}, (E_\alpha)_{\alpha \in Q_1})$ such that:
\begin{enumerate}[label=$\bullet$, itemsep=1mm]
    \item $E_q$ is a $\mathbb{K}$-vector space, for all $q \in Q_0$;
    \item  $E_\alpha : E_{s(\alpha)} \longrightarrow E_{t(\alpha)}$ is a $\mathbb{K}$-linear ma, for any $\alpha \in Q_1$; and,
    \item by setting $E_p = E_{\alpha_k} \circ \cdots \circ E_{\alpha_1}$ for any path $p = \alpha_k \ldots \alpha_1$ of $Q$, whenever we have a relation $k_1 p_1 + \ldots + k_s p_s$ on $Q$ in $R$, then $k_1 E_{p_1} + \ldots  + k_s E_{p_s} = 0$.
\end{enumerate}
This representation $E$ is \new{finite-dimensional} if, for all $q \in Q_0$, $E_q$ is finite-dimensional. Its \new{dimension vector} $\vdim(E)$ is the vector $(\dim(E_q))_{q \in Q_0} \in \mathbb{N}^{\#Q_0}$. Now, we assume that all the representations we consider are finite-dimensional. Given $E$ and $F$ two representations of $Q$, we define a \new{morphism} $\varphi : E \longrightarrow F$ as a collection $\varphi = (\varphi_q)_{q \in Q_0}$ for linear maps $\varphi_q : E_q \longrightarrow E_q$ such that, for all $\alpha \in Q_1$, we have $\varphi_{t(\alpha)} E_\alpha = F_\alpha \varphi_{s(\alpha)}$. We say that $E$ and $F$ are \new{isomorphic} whenever there exists a bijective morphism $\varphi : E \longrightarrow F$, or equivalently, if there exists a morphism $\varphi : E \longrightarrow F$ such that, for any $q \in Q_0$, $\varphi_q$ is an isomorphism of $\mathbb{K}$-vector spaces. In such a case, we write $E\cong F$. We obtain the structure of a category by endowing the collection of representations of $Q$ with the morphisms between them. In the following, we denote this category by $\rep(Q,R)$. In the case where $R = \varnothing$, we denote this category by $\rep(Q)$. Note that $\rep(Q, R)$ depends on the field $\mathbb{K}$.

 Any (basic) finite-dimensional $\mathbb{K}$-algebra is isomorphic to the quotient algebra $\mathbb{K}Q/\langle R \rangle$ for some bounded quiver $(Q,R)$ where $Q$ is finite, and $\langle R \rangle$ is admissible. Moreover, the category $\rep(Q,R)$ is equivalent to the category of finitely generated left module over $\mathbb{K}Q/\langle R \rangle$. In particular, it implies that $\rep(Q,R)$ is an abelian Krull--Schmidt category with enough projective objects. We refer the reader to \cite{ASS06} for more details. 
 
 Write $E \oplus F$ for the direct sum of $E$ and $F$ in $\rep(Q,R)$. We say that a nonzero representation $X \in \rep(Q,R)$ is \new{indecomposable} if whenever we have $X \cong E \oplus F$ with some $E,F \in \rep(Q,R)$, then $E \cong 0$ or $F \cong 0$. Write $\ind(Q,R)$ for the collection of isomorphism classes of indecomposable representations of $(Q,R)$. Whenever $R = \varnothing$, we denote this collection by $\ind(Q)$. A \new{representation-finite type} bounded quiver is a bounded quiver that has finitely many isomorphism classes of indecomposable representations. 

 For any $\pmb{d} = (d_q)_{q \in Q_0} \in \mathbb{N}^{\#Q_0}$, we consider the following affine space
 \[\rep((Q,R),\pmb{d}) = \prod_{\alpha \in Q_1} \Hom \left(\mathbb{K}^{d_{s(\alpha)}}, \mathbb{K}^{d_{t(\alpha)}} \right) = \prod_{\alpha \in Q_1}  \Mat_{d_{t(\alpha)} \times d_{s(\alpha)}} (\mathbb{K}),\] and refer to it as the \new{representation space} of $(Q,R)$ with dimension vector $\pmb{d}$. Indeed, by choosing a basis for each of the vector spaces, we can identify the representations of $(Q,R)$ with the points of $\rep((Q,R), \pmb{d})$. The isomorphism classes of representations with vector dimension $\pmb{d}$ correspond exactly to the orbits of the group action given by the algebraic group $\vGL_{\pmb{d}}(\mathbb{K}) = \prod_{q \in Q_0} \GL_{d_q}(\mathbb{K})$, which acts on $\rep((Q,R), \pmb{d})$ by changing bases at each vertex.

\subsection{ADE Dynkin type quivers}
\label{ss:Dynkin}

In this paper, we show examples of our results on representation-finite quivers. We shortly recall a complete characterization of them due to P. Gabriel \cite{G72}.

\begin{definition} \label{def:Dynkin}
    Let $n \in \mathbb{N}^*$. A quiver $Q$ is said to be:
    \begin{enumerate}[label=$\bullet$, itemsep=1mm]
        \item of \new{$A_n$ type} if the underlying graph of $Q$ is of the following shape.
         \begin{center}
		\begin{tikzpicture}[-,line width=0.5mm,>= angle 60,color=black,scale=0.8]
			\node (1) at (0,0){$1$};
			\node (2) at (2,0){$2$};
			\node (3) at (4,0){$\cdots$};
			\node (4) at (6,0){$n$};
			\draw (1) -- (2) -- (3) -- (4);
		\end{tikzpicture}
	\end{center}
        \item of \new{$D_n$ type}, for $n \geqslant 4$, if the underlying graph of $Q$ is of the following shape.
        \begin{center}
		\begin{tikzpicture}[-,line width=0.5mm,>= angle 60,color=black,scale=0.8]
			\node (1) at (0,0){$1$};
			\node (2) at (2,0){$2$};
			\node (3) at (4,0){$\cdots$};
			\node (4) at (6,0){$n-2$};
            \node (5) at (7.4142, 1.4142){$n-1$};
            \node (6) at (7.4142, -1.4142){$n$};
			\draw (1) -- (2) -- (3) -- (4) -- (5);
            \draw (4) -- (6);
		\end{tikzpicture}
	\end{center}
        \item of \new{$E_6$}, \new{$E_7$} or \new{$E_8$ type} if the shape of the underlying graph of $Q$ is one of the following.
        \begin{center}
        \scalebox{0.7}{\begin{tikzpicture}[-,line width=0.5mm,>= angle 60,color=black,scale=0.7]
 		\node (1) at (0,0){$1$};
 		\node (3b) at (4,2){$4$};
 		\node (2) at (2,0){$2$};
 		\node (3) at (4,0){$3$};
 		\node (4) at (6,0){$5$};
 		\node (5) at (8,0){$6$};
 		\draw (1) -- (2) -- (3) -- (4) -- (5);
 		\draw (3b) -- (3);
 	      \begin{scope}[xshift=10cm]
 		\node (1) at (0,0){$1$};
 		\node (3b) at (4,2){$4$};
 		\node (2) at (2,0){$2$};
 		\node (3) at (4,0){$3$};
 		\node (4) at (6,0){$5$};
 		\node (5) at (8,0){$6$};
 		\node (6) at (10,0){$7$};
 		\draw (1) -- (2) -- (3) -- (4) -- (5) -- (6);
 		\draw (3b) -- (3);
 	\end{scope}
            \begin{scope}[xshift = 5cm, yshift = -3cm]
 		\node (1) at (0,0){$1$};
 		\node (3b) at (4,2){$4$};
 		\node (2) at (2,0){$2$};
 		\node (3) at (4,0){$3$};
 		\node (4) at (6,0){$5$};
 		\node (5) at (8,0){$6$};
 		\node (6) at (10,0){$7$};
 		\node (7) at (12,0){$8$};
 		\draw (1) -- (2) -- (3) -- (4) -- (5) -- (6) -- (7);
 		\draw (3b) -- (3);
        \end{scope}
 	\end{tikzpicture}}
     \end{center}
    \end{enumerate}
We say that $Q$ is an \new{$ADE$ Dynkin type} quiver whenever $Q$ is one of the above types.
\end{definition}

We have the following main result.

\begin{theorem}[{\cite{G72}}] \label{thm:Gab} Let $Q$ be a quiver. Then $Q$ is a representation-finite type quiver if and only if $Q$ is an $ADE$ Dynkin type quiver.
\end{theorem}

\begin{remark}
    This result holds among bounded quivers $(Q,R)$ with $R = \varnothing$.
\end{remark}

In the following, we recall the precise description of $\rep(Q)$ for $Q$ a type $A$ quiver. The reader can find those results in \cite{S14}.
 
Let $n \in \mathbb{N}^*$. The \new{intervals} of $\{1, \ldots, n\}$ are the sets $\llrr{i,j} := \{i, i+1, \ldots,j \}$ given by all $1 \leqslant i \leqslant j \leqslant n$. If $i=j$, we write $\llrr{i,i} = \llrr{i}$. Denote by $\mathcal{I}_n$ the set of intervals in $\{1, \ldots n\}$. For $K = \llrr{i,j} \in \mathcal{I}_n$, set $b(K) = i$ and $e(K) = j$. Call \new{interval set} any subset of $\mathcal{I}_n$.
	\begin{definition}\label{def:intervalbotandtop}  Let $Q$ be an $A_n$ type quiver. Consider $K,L \in \mathcal{I}_n$ such that $K \subseteq L$. We say that:
		\begin{enumerate}[label = $\bullet$]
			\item $K$ is \new{above} $L$ (\new{relative to} $Q$) if the following two assertions are satisfied:\begin{enumerate}[label = $\bullet$]
				\item $b(K) = b(L)$ or we have the arrow $b(K)-1 \longleftarrow b(K)$ in $Q$;
				\item $e(K) = e(L)$ or we have the arrow $e(K) \longrightarrow e(K)+1$ in $Q$.
			\end{enumerate}
			\item $K$ is \new{below} $L$ (\new{relative to} $Q$) if the following two assertions are satisfied:\begin{enumerate}[label = $\bullet$]
				\item $b(K) = b(L)$ or we have the arrow $b(K)-1 \longrightarrow b(K)$ in $Q$;
				\item $e(K) =e(L)$ or we have the arrow $e(K) \longleftarrow e(K)+1$ in $Q$.
			\end{enumerate}
		\end{enumerate}
	\end{definition}
	\begin{ex} Consider the following quiver.
		\begin{center}
			\begin{tikzpicture}[->,line width=0.5mm,>= angle 60,color=black,scale=0.8]
				\node (Q) at (-1,0){$Q =$};
				\node (1) at (0,0){$1$};
				\node (2) at (2,0){$2$};
				\node (3) at (4,0){$3$};
				\draw (1) -- (2);
				\draw (2) -- (3);
			\end{tikzpicture}
		\end{center}
		Then $\llrr{2}$ is above $\llrr{2,3}$ and below $\llrr{1,2}$.
	\end{ex}
	Note that any interval is above and below itself, relative to all $A_n$ type quivers.
	Let $Q$ be a quiver of $A_n$ type. To any interval $K \in  \mathcal{I}_n$, we consider $X_K$ the representation of $Q$ defined as it follows:
	\begin{enumerate}[label = $\bullet$]
		\item $(X_K)_q = \mathbb{K}$ if $q \in K$, $(X_K)_q = 0$ otherwise;
		\item $(X_K)_{\alpha} = \Id_\mathbb{K}$ if $\alpha$ is such that $\{s(\alpha), t(\alpha)\} \subseteq K$, $(X_K)_{\alpha} = 0$ otherwise;
	\end{enumerate}
	Note that $X_K$ is an indecomposable representation of $Q$ for all $K \in \mathcal{I}_n$.
	\begin{theorem}\label{thm:indecandmorphtypeA} Let $Q$ be an $A_n$ type quiver. \begin{enumerate}[label = $(\alph*)$]
			\item \label{indectypeA} The isomorphism classes of indecomposable representations of $Q$ are in bijection with $\mathcal{I}_n$; more precisely, they are described by indecomposable representations $X_K$ for $K \in \mathcal{I}_n$;
			\item \label{morphtypeA} The homomorphism space between two indecomposable representations of $Q$ is of dimension at most one; more precisely, $\Hom(X_K, X_L)$ is nonzero if and only if  there exists an interval $J$ such that $J$ is above $K$ and below $L$ relative to $Q$; if such an interval exists, it is unique and 
			$\Hom(X_K,X_L)$ consists of scalar multiples of the morphism $\phi = (\phi_q)_{q \in Q_0}$ such that $\phi_q = \Id_\mathbb{K}$ if $q \in J$ and $\phi_q = 0$ otherwise. 
		\end{enumerate}
	\end{theorem}
	In the light of the previous result : \begin{enumerate}[label = $\bullet$]
		\item for all interval sets $\mathscr{J} \subseteq \mathcal{I}_n$ and all quivers $Q$ of $A_n$ type, write $\Cat_Q(\mathscr{J})$ for the subcategory of $\rep(Q)$ additively generated by $X_K$ for $K \in \mathscr{J}$;
		\item for all quivers $Q$ of $A_n$ type and for all nonzero subcategories $\mathscr{C}$ of $\rep(Q)$, let $\Int(\mathscr{C})$ be the interval set of $\mathcal{I}_n$ consisting of intervals $K$ such that $X_K \in \mathscr{C}$.
	\end{enumerate}
	Under \cref{conv:addsubcat}, the subcategories we are considering are additively generated by $X_K$ for $K \in \mathscr{J}$ for some  $\mathscr{J} \subset \mathcal{I}_n$.
	
	Hence, for any $A_n$ type quiver $Q$, for all $\mathscr{J} \subseteq \mathcal{I}_n$ and for all subcategories $\mathscr{C} \subseteq \rep(Q)$, we have $\mathscr{J} = \Int(\Cat_Q(\mathscr{J})) \text{ and } \mathscr{C} = \Cat_Q(\Int(\mathscr{C})).$

    Now we recall the description of all the short exact sequences between indecomposable representations.

    \begin{theorem} \label{thm:exttypeA} Let $Q$ be an $A_n$ type quiver. Let $D,F \in \ind(Q)$ and consider a non-split exact sequence \[\begin{tikzcd}
	0 & D & E & F &  0.
	\arrow[from=1-1, to=1-2]
	\arrow[tail, from=1-2, to=1-3]
	\arrow[two heads,from=1-3, to=1-4]
	\arrow[from=1-4, to=1-5]
\end{tikzcd} \]
      Then exactly one of the following assertions holds:
      \begin{enumerate}[label=$(\alph*)$,itemsep=1mm]
          \item We have $E \in \ind(Q)$, and if we consider $K,L \in \mathcal{I}_n$ such that $D \cong X_K$ and $F \cong X_L$, then $K \cup L \in \mathcal{I}_n$, $K \cap L = \varnothing$ and $E \cong X_{K \cup L}$; or,
          \item We have $E \cong E_1 \oplus E_2$ with $E_1,E_2 \in \ind(Q)$, and there exist $K,L \in \mathcal{I}_n$ with $K \cap L \neq \varnothing$ such that both assertions hold:
          \begin{enumerate}[label=$\bullet$,itemsep=1mm]
          \item $K,L,K \cap L$ and $K \cup L$ are four different intervals in $\mathcal{I}_n$; and,
          \item either $\{D,F\} = \{X_K, X_L\}$ and $\{E_1,E_2\} = \{X_{K\cap L}, X_{K \cup L} \}$, or \\
          $\{D,F\} = \{X_{K \cap L}, X_{K \cup L}\}$ and $\{E_1,E_2\} = \{X_{K}, X_{L} \}$.
          \end{enumerate}
          In this case, we call it a \new{diamond exact sequence.} 
      \end{enumerate}
    \end{theorem}

\begin{definition}[Diamond exact sequence]
Let $Q$ be an $A_n$ type quiver. A non-split short exact sequence
\[
0 \longrightarrow D \longrightarrow E \longrightarrow F \longrightarrow 0
\]
in $\rep(Q)$ is called a \emph{diamond exact sequence} if
\begin{enumerate}[label=(\roman*), itemsep=1mm]
    \item $D,F \in \ind(Q)$;
    \item $E$ is decomposable and,
    \[
    E \cong E_1 \oplus E_2 \quad \text{with } E_1,E_2 \in \ind(Q).
    \]
\end{enumerate}
\end{definition}
    
